\item \textbf{Kết quả không phải là một bản chụp tức thời của cả hệ thống.} \texttt{filecollect()} chạy trước \texttt{proccollect()}, và mỗi ô chỉ bị khóa trong lúc chép. Nếu có process khác mở, đóng file hay thoát cùng lúc, hai danh sách có thể không khớp nhau, ví dụ một file vẫn liệt kê pid của process vừa đóng nó (mục 7 có một trường hợp thật). Muốn chụp nhất quán thì phải giữ cùng lúc \texttt{ftable.lock} và khóa của mọi process, rất dễ deadlock. Nhóm chấp nhận điều này vì \texttt{sysinfo} chỉ dùng để quan sát.

\item \textbf{Đọc \texttt{ofile[]} của process khác.} \texttt{ofile[]} chỉ do chính process đó sửa và không có khóa riêng; giữ \texttt{p->lock} chỉ bảo đảm process không bị giải phóng trong lúc đọc. Con trỏ đọc được có thể vừa bị \texttt{close}, nhưng kernel chỉ so sánh con trỏ (\texttt{procs\_using\_file()}) hoặc tính chỉ số từ nó (\texttt{fileidx()}), không truy cập vào \texttt{struct file}, nên không có nguy cơ đọc vùng nhớ đã giải phóng.
